\answer{ex:02:1}{Ряд завдовжки $3.2$~мм із кроком $25\um$: $\approx129$ детекторів. Спрацьовує смуга завширшки $v\,\tau\ln2=1.7$~мм --- $\approx69$ детекторів, більше ніж половина ряду; напрямок --- центр смуги.}
\answer{ex:02:2}{Вікно $-\tau\ln\frac{w_2}{\theta-w_1}\le\delta\le\tau\ln\frac{w_1}{\theta-w_2}$, тобто $[-0.235,\,0.178]\ms$; центр $c=\frac\tau2\ln\frac{w_1(\theta-w_1)}{w_2(\theta-w_2)}=-28$~мкс. Напрямок зсувається на $28$~мкс у бік сильнішого вуха --- майже три кроки роздільності.}
\answer{ex:02:3}{Потрібні $s_i$ з $s_iw_{ij}s_j\ge0$; вони існують тоді й лише тоді, коли цикли парні. Взаємне гальмування зводиться (стійких коливань немає), петля \nt{GLUT}--\nt{GABA} --- ні (може коливатися).}
\answer{ex:02:4}{Шість нейронів $g_k=[x+y+z\ge k]$ і $\bar g_k=[\bar x+\bar y+\bar z\ge4-k]$, $k=1,2,3$; $C=g_2$, $\bar C=\bar g_2$, $S=[g_1+\bar g_2+g_3\ge2]$, $\bar S=[\bar g_1+g_2+\bar g_3\ge2]$: вісім нейронів. З І та АБО --- $12$.}
\answer{ex:02:5}{$T_{\min}(0.6)=0.718\,\tau$; $I_c=0.225$, біля порога $T_{\min}\propto(I-I_c)^{-1/2}$.}
\answer{ex:02:6}{(а)~$500$ ланок, $5\cdot10^4$ нейронів; розкид $4.5\ms$ ($0.45\,\%$), відносна похибка $\propto T^{-1/2}$. (б)~Спільний для всіх ланок випадковий множник затримки з розкидом $5\,\%$.}
\answer{ex:02:7}{Освіжати раз на $\tau_F\ln2=0.69\,\text{с}$: $4.3$ імпульсу за секунду на нейрон проти $20$, у $4.6$ раза ощадливіше ($4\cdot10^{-7}$ проти $2\cdot10^{-6}$~Дж при $10^{-10}$~Дж на імпульс). Тригер утрачає біт; конденсатор зберігає його, якщо глушіння почалося не пізніше ніж через $0.49\,\text{с}$ після освіження (імовірність $0.71$).}
